\chapter{Electroencephalogram (EEG)}
\index{Electroencephalogram (EEG)}

\section*{Definition and Overview}
An electroencephalogram, or EEG, is a test that records the electrical activity of the brain using electrodes placed on the scalp. The brain's billions of nerve cells produce tiny electrical signals as they communicate, and the EEG captures the overall pattern of this activity as wavy lines. The test is used mainly to investigate epilepsy and seizures, and also to help assess sleep disorders, confusion, brain injury, and other neurological conditions. An EEG is safe, painless, and does not involve radiation or electric currents entering the head.

\section*{Epidemiology}
The EEG is an important test in neurology, performed for people with suspected or known epilepsy, which affects about 1 in 100 people. It is also used in sleep laboratories and hospital neurology departments for the assessment of altered consciousness and brain injury. A significant proportion of people who have a seizure will have an EEG as part of their assessment, along with imaging and blood tests.

\section*{Aetiology and Pathophysiology}
The EEG reflects the summed electrical activity of the brain.
\begin{itemize}
    \item \textbf{Brain rhythms:} brain activity is described in frequency bands: alpha waves when awake and relaxed, beta waves during activity, and theta and delta waves during sleep and in some conditions
    \item \textbf{Normal patterns:} the EEG pattern changes with age, state of alertness, and mental activity
    \item \textbf{Epilepsy:} abnormal bursts of electrical discharge produce spikes and sharp waves characteristic of epilepsy
    \item \textbf{Other abnormalities:} slowing of rhythms occurs with brain injury, infection, metabolic disturbance, and drugs
    \item \textbf{Activation procedures:} hyperventilation, flashing lights, and sleep increase the chance of capturing abnormalities
    \item \textbf{Limitations:} many people with epilepsy have a normal EEG between seizures, and some EEG features are non-specific
\end{itemize}

\section*{Clinical Features}
An EEG is requested for specific clinical reasons.
\begin{itemize}
    \item Suspected epilepsy or a first seizure
    \item Classification of the epilepsy type to guide treatment
    \item Investigation of blackouts, funny turns, and unexplained episodes
    \item Assessment of confusion, delirium, or coma
    \item Evaluation of sleep disorders, including sleep-related seizures
    \item Monitoring during brain surgery and in intensive care
\end{itemize}

\section*{Differential Diagnosis}
The EEG is interpreted with the clinical picture.
\begin{itemize}
    \item Seizures arising from deep brain regions may not show on a scalp EEG
    \item Faints, panic attacks, and other events can mimic seizures
    \item Some normal variants can be mistaken for abnormality
    \item A normal EEG does not exclude epilepsy
\end{itemize}

\section*{When to Seek Medical Care}
An EEG is arranged by a specialist. Anyone who has had a possible seizure, an unexplained collapse, or prolonged confusion should be assessed by a doctor.

\textbf{Contact your local emergency services for:}
\begin{itemize}
    \item A seizure lasting more than five minutes or repeated seizures
    \item A first seizure in pregnancy or with difficulty breathing
    \item Prolonged unconsciousness
    \item Sudden weakness, facial drooping, or difficulty speaking
\end{itemize}

\section*{Diagnosis and Investigations}
The EEG procedure is straightforward.
\begin{itemize}
    \item \textbf{Preparation:} hair is washed and left free of products; a sleep-deprived EEG may be arranged to increase sensitivity
    \item \textbf{The test:} electrodes are attached to the scalp, and the recording is taken at rest, with eyes open and closed
    \item \textbf{Activation:} hyperventilation and photic stimulation are used
    \item \textbf{Duration:} usually 30 to 60 minutes
    \item \textbf{Longer monitoring:} sleep-deprived, 24-hour, and video-EEG recordings for harder-to-capture episodes
    \item \textbf{Reporting:} a clinical neurophysiologist or neurologist interprets the recording
\end{itemize}

\section*{Management}
The EEG informs diagnosis and treatment.
\begin{itemize}
    \item \textbf{Diagnosis of epilepsy:} characteristic discharges support the diagnosis
    \item \textbf{Classification:} the pattern distinguishes generalised from focal epilepsy
    \item \textbf{Treatment choice:} anti-epileptic medicines are selected partly on the EEG findings
    \item \textbf{Monitoring:} repeat recordings assess progress or clarify uncertain cases
    \item \textbf{Context:} the result is always weighed with the history and other tests
\end{itemize}

\section*{Complications}
\begin{itemize}
    \item The test is safe and painless
    \item Mild skin irritation from the electrodes in some people
    \item The risk of misinterpretation without specialist review
    \item The inconvenience of longer monitoring studies
\end{itemize}

\section*{Prognosis and Outcomes}
The EEG is a valuable, reliable tool in the assessment of epilepsy and other neurological conditions. When used and interpreted appropriately, it helps establish accurate diagnoses, guides effective treatment, and clarifies the nature of worrying episodes, contributing to better outcomes for people with neurological disorders.

\section*{Prevention and Self-Care}
\begin{itemize}
    \item Follow the preparation instructions, including sleep advice
    \item Bring a list of medicines and a description of any episodes
    \item Ask about the results and what they mean for your care
    \item Do not stop epilepsy medicines without medical advice
    \item Attend follow-up appointments as arranged
\end{itemize}

\section*{Sources and Further Reading}
The following reputable sources provide further information for healthcare students and patients seeking reliable, current advice about the EEG.
\begin{itemize}
    \item Epilepsy Australia. \textit{EEG testing information.} https://www.epilepsyaustralia.net/
    \item Healthdirect Australia. \textit{Electroencephalogram (EEG).} https://www.healthdirect.gov.au/eeg
    \item NHS. \textit{Electroencephalogram (EEG).} https://www.nhs.uk/conditions/electroencephalogram/
    \item Mayo Clinic. \textit{EEG (electroencephalogram).} https://www.mayoclinic.org/tests-procedures/eeg/
    \item MedlinePlus. \textit{Electroencephalogram.} https://medlineplus.gov/ency/article/003931.htm
    \item Better Health Channel. \textit{EEG brain function test.} https://www.betterhealth.vic.gov.au/
\end{itemize}
